% ============================================================
% SECTION 8: RENAISSANCE & EARLY MODERN (~1500 - 1700)
% Slides 56-65 (10 slides)
% Includes: Cardano, Viete, Napier, Descartes, Fermat,
%           Newton, Leibniz, Bernoullis, Calculus summary
% ============================================================

% ---- Section Divider ----
{
\setbeamercolor{background canvas}{bg=renaissanceGold}
\begin{frame}[plain]
  \centering
  \vfill
  \textcolor{white}{\Huge\bfseries The Modern Age Begins}\\[0.7em]
  \textcolor{white!85}{\Large Revolution in Symbol, Method, and Thought}\\[1.5em]
  \textcolor{white!65}{\large 1500--1700 CE}
  \vfill
  \visualplaceholder[5cm]{Diagram}
  \vfill
\end{frame}
}

\section{Renaissance \& Early Modern Mathematics}

% ---- Slide 57: Cardano and the Cubic ----
\begin{frame}[shrink=25]{Cardano and the Cubic Equation --- A Mathematical Drama}
  \begin{columns}[T]
    \column{0.50\textwidth}
    \begin{personbox}[Gerolamo Cardano]{renaissanceGold}
      {\small (Latin: Hieronymus Cardanus)}\\[3pt]
      \textbf{Dates:} 1501 -- 1576 CE\\
      \textbf{Origin:} Pavia, Italy
    \end{personbox}
    \vspace{0.3em}
    \begin{itemize}\setlength\itemsep{2pt}
      \item Published the general solution to \textbf{cubic equations} in \textit{Ars Magna} (1545)
      \item But who really discovered it?
      \item The solution forced the invention of \textbf{complex numbers} ($\sqrt{-1}$)
    \end{itemize}

    \vspace{0.3em}
    % Cubic formula (simplified form)
    \visualplaceholder[5cm]{Diagram}

    \column{0.47\textwidth}
    % Drama triangle TikZ
    \visualplaceholder[5cm]{Arrows with labels}
  \end{columns}

  \note{
    \textbf{Talking point:} ``The solution of the cubic equation is one of the great dramas
    in math history --- involving disputed promises, public math duels, and a formula so
    complicated it forced the invention of imaginary numbers.''

    Del Ferro discovered the solution first ($\sim$1515) but kept it secret.
    Tartaglia rediscovered it independently. Cardano obtained it from Tartaglia,
    then published after finding del Ferro's priority in his papers.

    \verifytag{The exact nature of Cardano's promise to Tartaglia is contested.
    Some sources say he swore an oath of secrecy; others say the commitment was more
    ambiguous. Cardano later argued that since del Ferro had discovered it first,
    the oath did not apply.}

    OPTIONAL INTERACTIVE: ``Quick question for the class: Cardano promised
    Tartaglia he would keep the cubic solution secret. Then he published it.
    Was Cardano wrong? Or was the math too important to keep secret?''
    [2--3 hands. Brief discussion. This breaks up Session 3's exposition.]
  }
\end{frame}

% ---- Slide 58: Viete ----
\begin{frame}[shrink=8]{Fran\c{c}ois Vi\`{e}te --- Symbolic Algebra}
  \begin{columns}[T]
    \column{0.50\textwidth}
    \begin{personbox}[Fran\c{c}ois Vi\`{e}te]{renaissanceGold}
      {\small (Latin: Franciscus Vieta)}\\[3pt]
      \textbf{Dates:} 1540 -- 1603 CE\\
      \textbf{Origin:} Fontenay-le-Comte, France
    \end{personbox}
    \vspace{0.3em}
    \begin{itemize}
      \item Introduced systematic use of \textbf{letters} for \emph{both} unknowns and known quantities
      \item Vowels for unknowns, consonants for knowns
      \item A major step toward \textbf{modern algebraic notation}
      \item Also worked on trigonometric identities and cryptanalysis
    \end{itemize}

    \column{0.47\textwidth}
    % Before and after comparison
    \visualplaceholder[5cm]{Medieval style}
  \end{columns}

  \note{
    \textbf{Talking point:} ``Before Vi\`{e}te, algebra was mostly done in words and sentences.
    After Vi\`{e}te, it started to look like what you see in your textbook.''

    \verifytag{Vi\`{e}te's innovation was the consistent use of letters for BOTH unknowns
    and known parameters, creating a general symbolic framework. Earlier uses (Diophantus,
    Islamic algebraists) were more limited. Qualify ``introduced'' accordingly.}
  }
\end{frame}

% ---- Slide 59: Napier ----
\begin{frame}[shrink=25]{John Napier --- Logarithms Change Everything}
  \begin{columns}[T]
    \column{0.52\textwidth}
    \begin{personbox}[John Napier]{renaissanceGold}
      \textbf{Dates:} 1550 -- 1617 CE\\
      \textbf{Origin:} Edinburgh, Scotland
    \end{personbox}
    \vspace{0.3em}
    \begin{itemize}
      \item Invented \textbf{logarithms} (published 1614)
      \item Transformed computation: turned \textbf{multiplication into addition}
      \item Laplace: logarithms ``doubled the life of astronomers'' by halving their computational labor
    \end{itemize}

    \vspace{0.3em}
    % Logarithm property
    \visualplaceholder[5cm]{Diagram}

    \column{0.45\textwidth}
    \visualplaceholder[4cm]{Page from Napier's\\logarithm tables (1614)\\
      \textit{Mirifici Logarithmorum\\Canonis Descriptio}}

    \vspace{0.5em}
    \visualplaceholder[4cm]{``Napier's bones'' ---\\his mechanical multiplication device}
  \end{columns}

  \note{
    \textbf{Talking point:} ``Before calculators, before computers --- there were logarithms.
    They were the computational revolution of the 17th century.''

    Logarithm tables were the essential computational tool for scientists, navigators,
    and engineers for over 300 years, until electronic calculators replaced them in the 1970s.
  }
\end{frame}

% ---- Slide 60: Descartes ----
\begin{frame}{Ren\'{e} Descartes --- Geometry Meets Algebra}
  \begin{columns}[T]
    \column{0.50\textwidth}
    \begin{personbox}[Ren\'{e} Descartes]{renaissanceGold}
      {\small (Latin: Renatus Cartesius)}\\[3pt]
      \textbf{Dates:} 1596 -- 1650 CE\\
      \textbf{Origin:} La Haye en Touraine, France;\\lived in the Dutch Republic
    \end{personbox}
    \vspace{0.3em}
    \begin{itemize}
      \item Created \textbf{analytic geometry} --- the merger of algebra and geometry
      \item Published in \textit{La G\'{e}om\'{e}trie} (1637)
      \item \textbf{Cartesian coordinates}: describe curves as equations, solve geometric problems algebraically
    \end{itemize}

    \column{0.47\textwidth}
    % Cartesian coordinate system with a curve (pure TikZ)
    \visualplaceholder[5cm]{Grid lines}
  \end{columns}

  \note{
    The merger of algebra and geometry was arguably the single most important conceptual
    advance since the Greeks. It made calculus possible.

    Portrait source: Frans Hals painting (Louvre, Wikimedia Commons).
  }
\end{frame}

% ---- Slide 61: Fermat ----
\begin{frame}[shrink=25]{Pierre de Fermat --- The Prince of Amateurs}
  \begin{columns}[T]
    \column{0.52\textwidth}
    \begin{personbox}[Pierre de Fermat]{renaissanceGold}
      \textbf{Dates:} 1601 -- 1665 CE\\
      \textbf{Origin:} Beaumont-de-Lomagne, France
    \end{personbox}
    \vspace{0.3em}
    \begin{itemize}
      \item Co-founded \textbf{probability theory} (with Pascal, 1654)
      \item Developed early \textbf{differential calculus} methods (tangents, maxima/minima)
      \item Stated \textbf{Fermat's Last Theorem} in a marginal note (c.\ 1637)
    \end{itemize}

    \vspace{0.3em}
    % The famous marginal note
    \visualplaceholder[5cm]{Diagram}

    \column{0.45\textwidth}
    % Fermat's Last Theorem visual
    \visualplaceholder[5cm]{Timeline}

    \vspace{0.5em}
    \visualplaceholder[3.5cm]{The Fermat--Pascal correspondence\\on probability (1654)}
  \end{columns}

  \note{
    \textbf{Talking point:} ``His margin note launched a 358-year quest that became the most
    famous unsolved problem in mathematics.''

    Fermat was a lawyer by profession --- mathematics was his hobby. He communicated most of
    his results in letters, not formal publications.
  }
\end{frame}

% ---- Slide 62: Newton ----
\begin{frame}[shrink=4]{Isaac Newton --- The Method of Fluxions}
  \begin{columns}[T]
    \column{0.52\textwidth}
    \begin{personbox}[Sir Isaac Newton]{renaissanceGold}
      \textbf{Dates:} 1642 -- 1726/27 CE\\
      \textbf{Origin:} Woolsthorpe-by-Colsterworth, England
    \end{personbox}
    \vspace{0.3em}
    \begin{itemize}
      \item Co-invented \textbf{calculus} (``the method of fluxions''), developed independently during 1665--1666
      \item Published \textit{Principia Mathematica} (1687) --- used calculus to describe gravity and planetary motion
      \item Also: optics, reflecting telescope
    \end{itemize}

    \column{0.45\textwidth}
    \visualplaceholder[4cm]{Page from the \textit{Principia}\\(Cambridge University Library)}

    \vspace{0.5em}
    % Newton's notation (using tcolorbox to avoid TikZ scope issue)
    \begin{tcolorbox}[colframe=renaissanceGold!60,colback=renaissanceGold!5,boxrule=0.5pt,arc=3pt]
      \centering
      \textbf{Newton's notation:}\\[4pt]
      $\dot{x}$ \quad (fluxion of $x$)\\[2pt]
      $\ddot{x}$ \quad (second fluxion)\\[6pt]
      {\footnotesize\textcolor{gray}{Still used in physics today for time derivatives}}
    \end{tcolorbox}
  \end{columns}

  \note{
    Newton developed calculus during the ``annus mirabilis'' (miracle year) of 1665--1666
    while Cambridge was closed due to plague. He was 23.

    The priority dispute with Leibniz became one of the ugliest feuds in the history of
    science, lasting decades and dividing British and Continental mathematics.
  }
\end{frame}

% ---- Slide 63: Leibniz ----
\begin{frame}[shrink=17]{Gottfried Wilhelm Leibniz --- The Notation that Won}
  \begin{columns}[T]
    \column{0.50\textwidth}
    \begin{personbox}[G.\ W.\ Leibniz]{renaissanceGold}
      \textbf{Dates:} 1646 -- 1716 CE\\
      \textbf{Origin:} Leipzig, Germany;\\worked in Hanover and Paris
    \end{personbox}
    \vspace{0.3em}
    \begin{itemize}
      \item Co-invented calculus \textbf{independently} of Newton (published 1684--86)
      \item His \textbf{notation} is what we use today
      \item Invented a mechanical calculator (Stepped Reckoner)
      \item Pioneered \textbf{binary arithmetic}
    \end{itemize}

    \column{0.47\textwidth}
    % Newton vs. Leibniz notation comparison
    \visualplaceholder[5cm]{Winner arrow}
  \end{columns}

  \vspace{0.2em}
  \begin{insightbox}
    {\footnotesize Every time you write $\frac{dy}{dx}$ or the integral sign $\int$, you are using \textbf{Leibniz's} language, not Newton's.}
  \end{insightbox}

  \note{
    \textbf{Talking point:} ``Newton may have gotten there first, but Leibniz's notation won.
    Every time you write $dy/dx$ or the integral sign, you are using Leibniz's language,
    not Newton's.''

    Leibniz also envisioned a ``universal language'' of reasoning --- a precursor to
    formal logic and computer science.
  }
\end{frame}

% ---- Slide 64: The Bernoulli Family ----
\begin{frame}[shrink=5]{The Bernoulli Family --- A Mathematical Dynasty}
  \begin{columns}[T]
    \column{0.45\textwidth}
    {\small The Bernoulli family of Basel, Switzerland: at least \textbf{eight} notable mathematicians across three generations.}

    \vspace{0.3em}
    \begin{personbox}[Jacob Bernoulli (1655--1705)]{renaissanceGold}
      {\footnotesize Basel, Switzerland\\
      Bernoulli numbers $\bullet$ \textit{Ars Conjectandi} (1713)\\
      The \textbf{law of large numbers}}
    \end{personbox}

    \vspace{0.2em}
    \begin{personbox}[Johann Bernoulli (1667--1748)]{renaissanceGold}
      {\footnotesize Basel, Switzerland\\
      \textbf{Taught Euler} $\bullet$ Brachistochrone problem\\
      Fierce rivalry with brother Jacob}
    \end{personbox}

    \column{0.52\textwidth}
    % Family tree
    \visualplaceholder[5cm]{Generation 1}

    \vspace{0.4em}
    % Brachistochrone curve
    \visualplaceholder[5cm]{start}
  \end{columns}

  \note{
    \textbf{Talking point:} ``The Bernoullis were the mathematical equivalent of a royal
    dynasty. Two brothers, Jacob and Johann, were bitter rivals --- and both were geniuses.
    Johann taught Euler, who became the most prolific mathematician in history.''

    The brachistochrone problem (1696): Johann posed the challenge ``find the curve of
    fastest descent.'' Newton famously solved it in a single evening.
  }
\end{frame}

% ---- Slide 65: Calculus Invented ----
\begin{frame}[shrink=12]{Calculus Invented --- The Floodgates Open}
  \centering

  {\small With calculus, analytic geometry, and modern notation in place,\\
  mathematics had the tools for \textbf{explosive growth}.}

  \vspace{0.6em}

  % Cascade timeline
  \visualplaceholder[5cm]{tiny 1637}

  \vspace{0.8em}

  \begin{columns}[T]
    \column{0.48\textwidth}
    {\small\textbf{What was now possible:}}
    \begin{itemize}\setlength\itemsep{2pt}
      \item[\textcolor{renaissanceGold}{\faCheck}] {\footnotesize Describe motion and change precisely}
      \item[\textcolor{renaissanceGold}{\faCheck}] {\footnotesize Compute areas, volumes, arc lengths}
      \item[\textcolor{renaissanceGold}{\faCheck}] {\footnotesize Model planetary orbits}
      \item[\textcolor{renaissanceGold}{\faCheck}] {\footnotesize Predict physical phenomena mathematically}
    \end{itemize}

    \column{0.48\textwidth}
    {\small\textbf{What came next:}}
    \begin{itemize}\setlength\itemsep{2pt}
      \item[\textcolor{enlightenmentTeal}{\faArrowRight}] {\footnotesize Euler: the most prolific mathematician ever}
      \item[\textcolor{enlightenmentTeal}{\faArrowRight}] {\footnotesize Lagrange \& Laplace: mechanics and celestial theory}
      \item[\textcolor{enlightenmentTeal}{\faArrowRight}] {\footnotesize Gauss: the ``Prince of Mathematicians''}
    \end{itemize}
  \end{columns}

  \vspace{0.5em}
  {\small\itshape ``With calculus invented, mathematics entered a new era of explosive growth\ldots''}

  \note{
    This is a transition slide. Emphasize: calculus was not invented by one person in
    one moment. It was a cascade of ideas spanning 50+ years, building on 2,000 years
    of Greek, Indian, and Islamic mathematics.

    The 18th century would see mathematics applied to every branch of science.
  }
\end{frame}
